\section{Introduction}
\label{sec:intro}

%-------------------------------------------------------------------------

Vision Transformers (ViTs)~\cite{dosovitskiy2021image} have emerged as a dominant architecture across computer vision, demonstrating strong performance in both recognition and dense prediction settings. However, ViTs suffer from quadratic computational complexity with respect to token count, making training at high resolutions prohibitively expensive. This limitation is particularly problematic for detection and segmentation, where fine-grained spatial detail is essential. Resolution extrapolation, i.e., training at low resolution and deploying at higher resolution, offers an appealing alternative, but is fundamentally constrained by the poor generalization of standard absolute positional embeddings beyond their training grid.

The practical importance of resolution extrapolation cannot be overstated. Object detection and semantic segmentation require precise localization of fine-grained structures, small objects, and thin boundaries, all of which benefit from higher-resolution input. Yet training a ViT-Base at $1024 \times 1024$ requires roughly $21\times$ the FLOPs compared to $224 \times 224$, making large-scale training at high resolution impractical for many research groups and deployment scenarios. If a model trained cheaply at low resolution could be deployed at higher resolution with minimal accuracy loss, it would dramatically reduce both the computational cost and the carbon footprint of modern vision systems. The critical bottleneck in realizing this goal is the positional encoding: while the feature extraction layers of a ViT are inherently resolution-agnostic (each patch is processed independently before attention), the positional encoding ties the model to a specific spatial grid and determines whether the learned spatial relationships generalize to unseen resolutions.

Recent progress in Large Language Models (LLMs)~\cite{touvron2023llama} has produced a rich landscape of positional encoding strategies specifically designed for length extrapolation. These methods fall into two broad families. \emph{Rotation-based methods}, such as RoPE~\cite{su2024roformer} and its extension YaRN~\cite{peng2024yarn}, encode relative positions by applying position-dependent complex rotations to query and key vectors, so that the attention score between two tokens depends only on their relative displacement. \emph{Additive bias methods}, such as ALiBi~\cite{press2022train} and FIRE~\cite{li2024functional}, instead inject positional information as a deterministic or learned bias added directly to attention logits, leaving the content-based query--key similarity untouched. Both families have enabled dramatic context-length extension in 1D language models, but their behavior on the two-dimensional spatial lattice of images, where relative offsets become 2D vectors and distances grow radially rather than linearly, remains largely unexplored.

A handful of prior works have adapted individual methods to vision, most notably 2D RoPE~\cite{heo2024rope}, but no study has systematically compared rotation-based and bias-based families under controlled conditions across multiple downstream tasks. This gap is significant because the structural properties that make a positional encoding extrapolatable in 1D (translation invariance, monotonic attention decay, suppression of high-frequency oscillations) may not straightforwardly transfer to 2D grids, where the topology, distance metric, and boundary conditions all differ.

In this work, we conduct a systematic empirical study of modern positional encoding strategies adapted to 2D Vision Transformers. Due to computational constraints, we focus on ViT-Tiny trained on ImageNet-100~\cite{deng2009imagenet} at $224 \times 224$, evaluating zero-shot extrapolation up to $1024 \times 1024$ ($4.57\times$ the training grid) without fine-tuning. We evaluate each method on classification, COCO object detection, and ADE20K semantic segmentation, and complement accuracy metrics with an attention-entropy analysis that reveals \emph{why} certain methods fail. Our contributions are:
\begin{itemize}
    \item We provide the first comprehensive comparison of RoPE~\cite{su2024roformer,heo2024rope}, ALiBi~\cite{press2022train}, YaRN~\cite{peng2024yarn}, and FIRE~\cite{li2024functional} for 2D resolution extrapolation in Vision Transformers, evaluating across classification, detection, and segmentation.
    \item We demonstrate that attention bias-based methods (ALiBi, FIRE) substantially outperform rotation-based methods (RoPE, YaRN) under extrapolation, with FIRE retaining $63.1\%$ accuracy at $4.6\times$ resolution versus RoPE's $18.9\%$.
    \item We analyze attention patterns to explain these differences, showing that bias-based methods maintain focused attention while rotation-based methods exhibit attention collapse at extrapolated resolutions.
    \item We identify an anomalous failure mode of ALiBi at very low resolutions (e.g., $6 \times 6$ patch grids), providing practical deployment guidance.
\end{itemize}
